\chapter{Development Method} \label{cap:method}

This work followed an iterative process organized into four phases: a study phase to understand GCC's
backend architecture, a requirements phase to define the instruction boundaries for each target, a
per-synthesis implementation-and-test cycle that forms the core of the work, and a validation phase
that confirms behavioral correctness on a reference simulator. Chapter~\ref{cap:requirements},
Chapter~\ref{cap:development}, and Chapter~\ref{cap:results} describe the outputs of these phases in
detail; this chapter describes the process itself.

\section{Study of GCC Internals}

The first phase was a structured study of GCC's backend extension mechanisms. The entry point was the
GCC Internals manual \cite{gcc-internals}, which documents the machine description language
(\texttt{.md}), the option description system (\texttt{.opt}), and the role of per-target
configuration headers. After establishing a conceptual model from the manual, the existing RISC-V
backend (\texttt{gcc/config/riscv/riscv.md}, \texttt{riscv.opt}, and \texttt{config.gcc}) was
examined as a working reference implementation.

Two diagnostic tools were central throughout the study and subsequent implementation phases. The
\texttt{-S} flag causes GCC to emit assembly rather than an object file, making the compiler's
instruction selection directly observable. The \texttt{-fdump-rtl-*} family of flags produces
snapshots of GCC's internal Register Transfer Language (RTL) representation at each compilation
stage; these were used to understand why a pattern matched or failed to match during instruction
selection, and to verify that \texttt{define\_expand} bodies fired at the correct point in the
compilation pipeline.

\section{Requirements Specification}

Requirements for each target were derived from two sources. For rvsc0 and rvsc1, the instruction
boundaries are externally fixed by the Hennessy-Patterson textbook \cite{patterson2020} and the
minimum ABI requirements, respectively. For rvsc2--rvsc7, the boundaries follow the natural RISC-V
extension hierarchy \cite{riscv-spec}. The complete requirements are specified in
Chapter~\ref{cap:requirements}.

\section{Synthesis Derivation and Implementation}

The rvsc0 and rvsc1 targets require GCC to synthesize every instruction outside the native set using
only the instructions the processor supports. Each synthesis was developed through the following cycle,
applied independently to each missing operation.

The process begins by identifying the missing operation, an instruction that GCC's middle-end may
legally request but that the target processor does not support. The synthesis algorithm is then derived
and expressed as C pseudocode, which serves both as a correctness argument and as an unambiguous
specification of the target behavior. The pseudocode is then lifted directly into a
\texttt{define\_expand} body in \texttt{riscv.md}, using GCC's RTL emit helpers to generate the
equivalent sequence of native instructions. Correctness of the expansion is confirmed in two steps:
the \texttt{-S} output is inspected to verify that no forbidden mnemonics appear, and the resulting
program is executed on the Spike RISC-V ISA simulator \cite{spike} to confirm that the synthesized
sequence produces the same result as the original instruction would have.

This cycle was repeated for each operation that the compiler may emit for a bare-metal freestanding C
program targeting rvsc0 or rvsc1. The derivations and the resulting implementation are described in
Chapter~\ref{cap:development}.

\section{Validation}

Validation addresses two correctness requirements defined in Chapter~\ref{cap:requirements}: ISA
compliance and behavioral equivalence.

ISA compliance is verified structurally by disassembling the compiled output and checking every
mnemonic against a per-target allowlist.

Behavioral equivalence is verified by differential execution: the same C program is compiled by both
the custom target and a reference RV32I compiler, and both binaries are executed on the Spike ISA
simulator \cite{spike}; the test passes when both produce identical exit codes.

A third validation layer exercises the compiler against a much larger, automatically generated corpus.
The GCC test suite includes a \textit{torture test} mode that systematically varies optimization flags
and source patterns across hundreds of programs; besides compiling \texttt{libgcc} against the custom
target also stresses constant materialization, multi-word arithmetic, and calling-convention edge cases
that hand-written behavioral tests would not cover.

Test results are presented in Chapter~\ref{cap:results}.
